\section*{Latihan dan Refleksi}
\addcontentsline{toc}{section}{Latihan dan Refleksi}

\subsection*{Latihan}
\begin{enumerate}[label=\arabic*.,leftmargin=*]
  \item Jelaskan perbedaan tujuan buku ajar ini dengan buku teks
        kompilator pada umumnya.
  \item Sebutkan empat CPMK dan hubungkan masing-masing dengan minimal
        satu bab inti.
  \item Mengapa buku ini memisahkan bab orientasi dan landasan teori
        dari empat belas bab inti?
  \item Gambarkan alur dari program sumber sampai AST dalam bentuk
        diagram sendiri.
  \item Baca \file{code/subsetc/examples/hello.sc}, lalu tuliskan
        token-token yang dihasilkan secara manual.
\end{enumerate}

\subsection*{Refleksi}
\begin{reflectionbox}
Jawab dengan jujur:
\begin{enumerate}[label=\arabic*.,leftmargin=*]
  \item Kompetensi apa yang paling ingin Anda capai pada mata kuliah ini?
  \item Bagian mana dari proses kompilasi yang paling sulit Anda bayangkan?
  \item Kebiasaan belajar apa yang akan Anda ubah setelah membaca bab ini?
\end{enumerate}
\end{reflectionbox}
